\section*{Chapitre \ref{ch:b1:elimination} --- La \texorpdfstring{$\beta$}{Beta}-élimination}

\begin{solution}{exo:b1:elimination:1}
2-Bromobutane~: but-1-ène, \cip{E}- et \cip{Z}-but-2-ène~; majoritaire~:
\cip{E}-but-2-ène. 2-Bromo-2-méthylbutane~: 2-méthylbut-2-ène
(majoritaire, trisubstitué) et 2-méthylbut-1-ène. Bromocyclohexane~:
cyclohexène seulement.
\end{solution}

\begin{solution}{exo:b1:elimination:2}
E1~: $v = k[\mathrm{RBr}]$~; E2~: $v = k[\mathrm{RBr}][\ce{EtO-}]$.
Mesurer la \omterm{def:b1:rate-laws:initial-rate}{vitesse initiale} de formation de l'alcène pour plusieurs
concentrations d'éthanolate~: elle croît proportionnellement pour l'E2,
elle ne change pas pour l'E1.
\end{solution}

\begin{solution}{exo:b1:elimination:3}
Le long de C2--C3, le Br porté par C2 peut être anti par rapport à l'un
ou l'autre hydrogène de C3 ou au méthyle C4~: les deux \omterm{def:b1:stereochemistry-in-depth:isomers}{conformations} où
un H est anti par rapport à Br peuvent réagir par E2 (en donnant les
alcènes \cip{E} et \cip{Z})~; la troisième ne peut pas éliminer vers
C3.
\end{solution}

\begin{solution}{exo:b1:elimination:4}
SN2 (primaire, hydroxyde, basse température)~; E2 (tertiaire, \omterm{def:b1:predominant-reaction:strong-acid}{base forte}
et encombrée)~; SN2 (iodure, bon \omterm{def:b1:electronic-effects:nucleophile}{nucléophile} et \omterm{def:b1:predominant-reaction:strong-acid}{base faible}, \omterm{def:b1:intermolecular-forces-solvents:solvent-classes}{solvant
aprotique})~; SN1 et E1, l'E1 augmentant avec le chauffage.
\end{solution}

\begin{solution}{exo:b1:elimination:5}
C3 ne porte qu'un hydrogène. Dans chaque diastéréoisomère, placer cet H
en anti par rapport au Br de C2 fixe les positions des autres groupes~:
le méthyle de C2 et l'éthyle de C3 sont du même côté dans un
diastéréoisomère et de part et d'autre dans l'autre. L'un donne le
\cip{E}-, l'autre le \cip{Z}-3-méthylpent-2-ène~: \omterm{def:b1:nucleophilic-substitution:stereospecific}{réaction
stéréospécifique}.
\end{solution}

\begin{solution}{exo:b1:elimination:6}
Le OH est protoné, \ce{R-OH2+}~; l'eau part, ce qui donne le
\omterm{def:b1:electronic-effects:intermediates}{carbocation} tertiaire \ce{(CH3)2C+-CH2CH3}~; une base (l'eau, \ce{HSO4-})
arrache un proton d'un carbone voisin. Produit majoritaire (Zaïtsev)~:
le 2-méthylbut-2-ène.
\end{solution}

\begin{solution}{exo:b1:elimination:7}
Le groupe \textit{tert}-butyle bloque la \omterm{def:b1:stereochemistry-in-depth:conformations}{chaise} où il est \omterm{def:b1:stereochemistry-in-depth:conformations}{équatorial}.
Dans l'isomère \textit{trans}, le brome est alors \omterm{def:b1:stereochemistry-in-depth:conformations}{équatorial} lui aussi,
jamais anti par rapport à une liaison C--H~: l'E2 doit attendre la
\omterm{def:b1:stereochemistry-in-depth:conformations}{chaise} inversée, très défavorable. Dans l'isomère \textit{cis}, le brome
est \omterm{def:b1:stereochemistry-in-depth:conformations}{axial}, avec deux hydrogènes \textit{trans}-diaxiaux~: E2 rapide.
\end{solution}

\begin{solution}{exo:b1:elimination:8}
L'éthanolate, petit, arrache l'hydrogène qui donne l'alcène le plus
substitué (Zaïtsev). Le \textit{tert}-butanolate, encombré, atteint plus
facilement les hydrogènes exposés du groupe méthyle que l'hydrogène du
\ce{CH2} interne~: l'alcène le moins substitué.
\end{solution}

\begin{solution}{exo:b1:elimination:9}
Le bromométhane n'a pas de carbone $\beta$. Dans le
1-bromo-2,2-diméthylpropane, le carbone $\beta$ est quaternaire et ne
porte aucun hydrogène.
\end{solution}

\begin{solution}{exo:b1:elimination:10}
Les deux produits exigent le \omterm{def:b1:electronic-effects:intermediates}{carbocation}, formé au cours de l'étape
lente à la vitesse
$k_1[\mathrm{RBr}]$. Il se partage ensuite entre la capture par
l'éthanol (\omterm{def:b1:rate-laws:order}{constante de vitesse} $k_S$) et la perte d'un proton (\omterm{def:b1:rate-laws:order}{constante
de vitesse} $k_E$), toutes deux d'\omterm{def:b1:rate-laws:apparent-order}{ordre apparent} 1 par rapport au
solvant~: la fraction d'alcène vaut $k_E/(k_E + k_S)$, indépendamment de
la concentration du \omterm{def:b1:nucleophilic-substitution:substitution}{substrat}.
\end{solution}

\begin{solution}{exo:b1:elimination:11}
À partir de la \omterm{def:b1:stereochemistry-in-depth:isomers}{conformation} du \omterm{def:b1:stereochemistry-in-depth:enantiomers}{composé méso} où les deux Br sont anti
(les deux phényles sont alors anti, et les deux H aussi), on tourne C1
de $120^\circ$ pour placer son H en anti par rapport au Br de C2. Les
deux groupes phényle sont alors du même côté de l'axe H--Br~: ils se
retrouvent \textit{cis}. Sur C1, Br l'emporte sur Ph~; sur C2, Ph
l'emporte sur H~; Br et le phényle de C2 sont \textit{trans}~:
\cip{E}-1-bromo-1,2-diphényléthène.
\end{solution}

\begin{solution}{exo:b1:elimination:12}
Seule la \omterm{def:b1:stereochemistry-in-depth:conformations}{chaise} où le \omterm{def:b1:electronic-effects:nucleophile}{groupe partant} est \omterm{def:b1:stereochemistry-in-depth:conformations}{axial} réagit~:
$v = k\,x\,[\mathrm{RY}][\mathrm{B}]$, où $x$ est la fraction de cette
\omterm{def:b1:stereochemistry-in-depth:conformations}{chaise}. Chaque groupe alkyle qu'elle oblige à être \omterm{def:b1:stereochemistry-in-depth:conformations}{axial} coûte plusieurs
kJ/mol (\qty{5.7}{kJ/mol} pour un méthyle), ce qui divise $x$ par dix
environ par groupe à température ambiante~: avec deux ou trois tels
groupes, $x$ tombe à $10^{-2}$ ou $10^{-3}$ et la réaction est d'autant
plus lente.
\end{solution}

\begin{solution}{pb:b1:elimination:1}
\textbf{1.} Une \omterm{def:b1:stereochemistry-in-depth:conformations}{chaise} où Cl, l'isopropyle et le méthyle sont tous
\omterm{def:b1:stereochemistry-in-depth:conformations}{équatoriaux}.
\textbf{2.} L'isopropyle et le méthyle \omterm{def:b1:stereochemistry-in-depth:conformations}{équatoriaux}, Cl \omterm{def:b1:stereochemistry-in-depth:conformations}{axial}.
\textbf{3.} Des \omterm{def:b1:stereochemistry-in-depth:enantiomers}{diastéréoisomères}~: ils ne diffèrent qu'en C1.
\textbf{4.} Le \omterm{def:b1:electronic-effects:nucleophile}{groupe partant} et l'hydrogène en $\beta$
\textit{trans}-diaxiaux.
\textbf{5.} Seulement dans la \omterm{def:b1:stereochemistry-in-depth:conformations}{chaise} inversée, où Cl, l'isopropyle et le
méthyle sont tous \omterm{def:b1:stereochemistry-in-depth:conformations}{axiaux}.
\textbf{6.} Les trois groupes \omterm{def:b1:stereochemistry-in-depth:conformations}{axiaux}~: le méthyle seul coûte environ
\qty{5.7}{kJ/mol}, le groupe isopropyle, plus gros, davantage, le chlore
un peu~: au total quelque \qty{13}{kJ/mol}, soit une fraction d'environ
$\exp(-13000/2479) = \num{5e-3}$, la valeur que retient le problème.
\textbf{7.} C2 (un H, il porte le groupe isopropyle) et C6 (deux H).
\textbf{8.} Cl étant \omterm{def:b1:stereochemistry-in-depth:conformations}{axial} sur C1, le H \omterm{def:b1:stereochemistry-in-depth:conformations}{axial} de C2 (l'isopropyle étant
\omterm{def:b1:stereochemistry-in-depth:conformations}{équatorial}) et le H \omterm{def:b1:stereochemistry-in-depth:conformations}{axial} de C6~: deux.
\textbf{9.} Dans la \omterm{def:b1:stereochemistry-in-depth:conformations}{chaise} triaxiale, le groupe isopropyle occupe la
position axiale de C2, donc son H est \omterm{def:b1:stereochemistry-in-depth:conformations}{équatorial}~: pas anti. C6 possède
un H \omterm{def:b1:stereochemistry-in-depth:conformations}{axial}~: anti.
\textbf{10.} Le chlorure de néomenthyle deux, le chlorure de menthyle
un.
\textbf{11.} Le long de C1--C6~: Cl (\omterm{def:b1:stereochemistry-in-depth:conformations}{axial}, vers le haut) sur C1 à
l'opposé du H \omterm{def:b1:stereochemistry-in-depth:conformations}{axial} (vers le bas) de C6~; les liaisons du cycle et le H
\omterm{def:b1:stereochemistry-in-depth:conformations}{équatorial} décalés entre eux.
\textbf{12.} Sur une \omterm{def:b1:stereochemistry-in-depth:conformations}{chaise}, aucune liaison C--H n'est \omterm{def:b1:elimination:periplanar}{synpériplanaire} à
la liaison C--Cl~; l'imposer exigerait un cycle éclipsé et tendu.
\textbf{13.} Arrachement du H de C2~: 4-méthyl-1-(propan-2-yl)cyclohex-1-ène
(trisubstitué)~; arrachement d'un H de C6~:
3-méthyl-6-(propan-2-yl)cyclohex-1-ène (disubstitué).
\textbf{14.} L'alcène trisubstitué, le plus substitué et le plus stable.
\textbf{15.} Les deux alcènes, dont \qty{78}{\percent} de
trisubstitué~: c'est cohérent.
\textbf{16.} Seulement l'alcène disubstitué~: l'inverse de la prévision
de Zaïtsev, imposé par le seul hydrogène anti disponible.
\textbf{17.} Les deux sont stéréospécifiques (exigence anti)~; le
chlorure de néomenthyle réagit de façon régiosélective (78/22)~; le
chlorure de menthyle donne un seul isomère de constitution.
\textbf{18.} Le carbone qui porte Cl est secondaire et flanqué d'un
groupe isopropyle~: il est encombré pour une SN2, alors que l'éthanolate
est une \omterm{def:b1:predominant-reaction:strong-acid}{base forte}.
\textbf{19.} $v \propto x\,n_{\mathrm{H}}$~: fraction de \omterm{def:b1:stereochemistry-in-depth:conformations}{chaise} réactive
multipliée par le nombre d'hydrogènes anti.
\textbf{20.} $0.95 \times 2 = 1.90$.
\textbf{21.} $\num{5.0e-3} \times 1 = \num{5.0e-3}$.
\textbf{22.} \textbf{$k(\text{néomenthyle})/k(\text{menthyle}) = 1.90/\num{5.0e-3}
= 380$}.
\end{solution}
